\documentclass[compress]{beamer}


\usetheme{Madrid}

\usepackage[english]{babel}
\usepackage[latin1]{inputenc}
\usepackage{times}
\usepackage[T1]{fontenc}
\usepackage{verbatim}
\usepackage{empheq}
\usepackage{cancel}

\usepackage{amsmath}

\DeclareMathOperator{\cm}{cm}
\DeclareMathOperator*{\Avg}{\mathbb{E}}
\newcommand{\dist}{\mathrm{dist}}
\newcommand{\diam}{\mathrm{diam}}
\newcommand{\oz}{{\bigcirc\!\!\!\!\mathrm{z}\,\,}}

% Author, Title, etc.

\newcommand{\mcite}[1]{{\color{blue}[#1]}}
\newcommand{\e}{\varepsilon}
\newtheorem{question}{Question}
\newtheorem{proposition}{Proposition}
\newtheorem{conjecture}{Conjecture}
\newtheorem{claim}{Claim}
\newtheorem{observation}{Observation}



\title[Non-linear spectral gaps]{Towards a calculus for non-linear spectral gaps}

\author[M.~Mendel]{Manor~Mendel, Open University}
\institute[OUI]{Jointly with Assaf Naor, NYU}
%\institute[OpenU, NYU]
%{
%  \inst{1}%
%  Open University of Israel
%  \and
%  \vskip-2mm
%  \inst{2}%
%  Courant Institute
%  }

\date[Nov '08]{Theory Seminar, HU, Nov '08}


\AtBeginSection
{
\begin{frame} %<beamer> \frametitle
{Outline}
\tableofcontents[currentsection]
\end{frame}
}
% The main document

\begin{document}

\begin{comment}
Title: Towards a calculus for non-linear spectral gaps
Speaker: Manor Mendel, The Open University of Israel

Abstract:
For a regular undirected graph G=(V,E), let g(G) be the largest positive number g
such that for every f:V--> Reals,

AVERAGE_{x,y \in V} |f(x)-f(y)|^2 < g^{-1} AVERAGE_{(x,y)\in E} |f(x) - f(y)|^2.

It is known that g(G) is the (normalized) spectral gap of G, i.e. g(G)=1 - \lambda_2(G), 
where \lambda_2(G) is the normalized second eigenvalue of G.

Motivated by questions on the embedability of expander graphs in more general metric spaces,
we replace the Reals in the above definition with other metric spaces, and study the behavior
of the resulting "non-linear spectral gap" under various operations.

We obtain the following results:
* A clean and general analysis of the zigzag product of graphs, devoid of linear algebra.
* A unification of geodesic spaces with non-positive curvature and uniformly convex Banach spaces as "barycentric" spaces.
* A proof that spectral gap in barycentric spaces of a power of graph increases similarly to the classical spectral gap.
* A high degree, yet non-trivial, family of graphs with constant spectral gap in K-convex Banach spaces. 
   The graph is obtained by applying the noise (Bonami-Beckner) operator on the Hamming cube, 
   and taking a quotient of it by the dual of "good" linear code.
   The analysis is inspired by a construction of Khot and Naor and the proof of Pisier's K-convexity theorem.

The above results are applied to show a construction of constant degree zigzag expanders in uniformly convex Banach spaces. 
As a side benefit, we also obtain a Lipschitz extension theorem from subsets of L_p or series/parallel graphs into spaces of non-positive curvature.

No prior knowledge of the notions above will be assumed.

Joint work with Assaf Naor (NYU).

\end{comment}


\begin{frame}
  \titlepage
\end{frame}

\begin{frame}{Outline}
  \tableofcontents
\end{frame}


\section{Expanders}

\begin{frame}{Graphs}

\begin{minipage}{2in}
$(n,d)$ graph $G=(V,E)$ is
\begin{itemize} 
\item on $n$ vertices. 
\item the degree is $d$. 
\item undirected, 
\item connected,
\item multigraph (self loops + parallel edges)
\end{itemize}
\end{minipage}\quad\qquad
\begin{minipage}{2in}
\includegraphics[width=2in]{graph-ex}
\end{minipage}
\medskip

The normalized adjacency matrix $A=(a_{ij})$
\[a_{ij}= \#((i,j)\in E) /d, \]
$A$ is symmetric doubly stochastic. Its eigenvalues are 
\[ 1=\lambda_1> \lambda_2 \ge \cdots \lambda_n \ge -1 .\]

Also define 
\( \lambda=\max\{\lambda_2, -\lambda_n\}=\max _{i\ge 2} |\lambda_i| .\)
\end{frame}


\begin{frame}{Spectral Gap}
\onslide<+->
\begin{definition}[Spectral gap]
\[ g(G)=1-\lambda_2(G) % \qquad \qquad  \tilde{g}(G)=1-\lambda(G) 
\]
\end{definition}

Intuition: 
$g(G)>0$ $\Longleftrightarrow$ $G$ is connected.

% $\tilde{g}(G)>0$ $\Longleftrightarrow$ $G$ is connected.
\onslide<+->
\begin{definition}[Expanders]
 An infinite family of graphs with uniform positive lower bound on the spectral gap.
\end{definition}

In this definition complete graphs are expanders.
We will usually be interested in \alert{constant degree} expanders.
Their existence is by no means trivial. However

\onslide<+->
\begin{theorem}[ Pinsker, Margulis]
Constant degree expanders do exist.
\end{theorem}
\end{frame}

\begin{frame}{Expanders are Common}
\begin{itemize}[<+->]
\item Derandomization
\item Networking 
\item Coding
\item Metric Geometry
\item Geometric Group Theory
\end{itemize}
\end{frame}

\begin{frame}{Poincar\'e Inequality (PI) for Expanders}
\begin{proposition}
Given $G=(V,E)$ and every $f:V\to \mathbb R$,
\[ \Avg_{x,y\in V} |f(x)-f(y)|^2  \le (1-\lambda_2)^{-1} \Avg_{(x,y)\in E} |f(x)-f(y)|^2. 
\]
\end{proposition}

\begin{overprint}
\begin{proof}[Proof (I\onslide*<2>{I})]
\onslide*<1>{
By translation, we assume wlog  that $\sum_x f(x)=0$, 
so 
\[ \Avg_{x,y\in V} |f(x)-f(y)|^2 =  2 \Avg_{x\in V} f(x)^2. \]

On the other hand,
\[
 \frac{1}{n} \sum_{(x,y)\in V} a_{xy} (f(x)-f(y))^2 
=\frac{2}{n} \sum_{x\in V} f(x)^2 - \frac{2}{n} \sum_{x,y\in V} a_{xy}f(x)f(y) .
\]
\renewcommand{\qed}{}}
\onslide*<2>{
Using the spectral decomposition $f=\sum_i \alpha_i v_i$, (remembering that $\alpha_1=\langle f, v_1 \rangle=0$), 
we have  
\begin{multline*} 
\sum_{x,y\in V} a_{xy}f(x)f(y)=
\sum_{x,y\in V} a_{xy} (\sum_i \alpha_i v_i(x)) (\sum_j \alpha_j v_j(y)) \\
= \sum_{i,j\ge 2} \alpha_i \alpha_j \sum_{x\in V}  v_i(x) \lambda_j v_j(x) 
= \sum_{i\ge 2}\alpha_i^2  \lambda_i \sum_{x\in V} v_i(x)^2
\le \lambda_2 \sum_{x\in V} f(x)^2.
\end{multline*}}
\end{proof}
\end{overprint}
\end{frame}

\begin{frame}{Redefining $g(G)$: PI Const}
The previous analysis is tight: Checking with $f=v_2$.

\onslide<+->
We can redefine
\[ g(G)^{-1} = \sup_{\text{const} \ne f:V\to \mathbb R } \frac{\Avg_{x,y\in V} |f(x)-f(y)|^2 }{\Avg_{(x,y)\in E} |f(x)-f(y)|^2}. \]

\onslide<+->
Furthermore
\[ g(G)^{-1} = \sup_{\text{const} \ne f:V\to  L_2 } \frac{\Avg_{x,y\in V} \|f(x)-f(y)\|_2^2 }{\Avg_{(x,y)\in E} \|f(x)-f(y)\|_2^2}. \]
\end{frame}

\section{Background \& Motivation}

\begin{frame}{BiLipschitz \& Coarse embedding}
\onslide<+->
Let $f:(X,\rho)\to (Y,\nu)$, satisfying
\begin{itemize}
\item  $\nu(f(x),f(y)) \ge \rho(x,y)$, $\forall x,y\in X$,
\item $\nu(f(x_0),f(y_0))=\rho(x_0,y_0)$ for some $x_0,y_0\in X$.
\end{itemize}

\begin{definition}[Distortion]
$\dist(f)= \sup_{x\ne y} \frac {\nu(f(x),f(y))}{\rho(x,y)} \qquad$
$ c_Y(X)= \inf \{\dist(f): f:X\to Y\}$.
\end{definition}

\medskip

\onslide<+->
\begin{definition}[Coarse Embedding]
$f:(X,\rho) \to (Y,\nu)$ is a coarse embedding if $\exists \alpha,\beta:[0,\infty) \to [0,\infty)$ such that
$\lim _{t\to \infty} \alpha(t)=\infty$, and
\[ \forall x,y\in X,\  \alpha(\rho(x,y)) \le \nu(f(x),f(y)) \le \beta(\rho(x,y)). \]
\end{definition}

\end{frame}

\begin{frame}{Non embedability of Expanders}
\onslide<+->
\begin{claim}
In $(n,d)$ graphs,  most pair of points are at distance $\Omega(\log_d n)$ (counting argument).
Therefore,
\[ \frac{\Avg_{x,y\in V} d_G(x,y)^2 }{\Avg_{(x,y)\in E} d_G(x,y)^2} = \Omega(\log_d ^2 n) .\]
\end{claim}

\onslide<+->
\begin{theorem}[Linial-London-Rabinoovich, Matou\v{s}ek, Gromov]
If $(G_n)_n$ is a  family of $(n,d)$ graphs  with uniform bound $g^{-1}$ on the PI constant. Then
\begin{itemize}
\item $c_{L_2}(G_n)=\Omega( \sqrt g \log_d n)$.
\item $(G_n)$ is not coarsely embeddable in $L_2$.
\end{itemize}
\end{theorem}

\onslide<+->
Non-embedding in other spaces?
\onslide <+-> Generalizes PI to those spaces!
\end{frame}



\begin{frame}{PI Constant in Abstract Spaces}
\onslide<+->
$\rho:X\times X \to  [0,\infty)$ is called symmetric kernel if $\rho(x,y)=\rho(y,x)$.

\begin{definition}[PI Constant]
\[ g(G,\rho)^{-1} = \sup_{\text{const} \ne f:V\to  X } \frac{\Avg_{x,y\in V} \rho(f(x),f(y)) }{\Avg_{(x,y)\in E} \rho(f(x),f(y))}. \]
\end{definition}

We will care mostly when $\rho$ is a power of a metric.

But, there are non trivial facts even in this abstract setting.

\onslide<+->

\begin{example}
\begin{itemize}
\item When $(X,\rho)=(\mathbb R, | \cdot |^2)$ or $(X,\rho)=(L_2, \| \cdot \|_2^2)$,
$g(G,\rho)$ is the spectral gap.
\item  When $(X,\rho)=(\{0,1\})$ or $(X,\rho)=(L_1, \|\cdot \|_1)$, 
$g(G,\rho)$ is (closely related to) the edge expansion. 
\end{itemize}
\end{example}
\end{frame}

\begin{frame}{Abstract PI: Simple Observations}
\onslide<+->
\begin{observation}[Cheeger-Tanner-Alon-Milman]
For any graph $G$ and non trivial $(X,\rho)$,  
\[ \sqrt{ g^{-1}(G)} \le g^{-1}(G,\{0,1\})\le g^{-1}(G,\rho). \]
\end{observation}

\onslide<+->
\begin{observation}
When $(X,\rho)$ is a metric space, and $p\ge 1$,
\[ g^{-1}(G,\rho^p)\le \diam(G)^p \]
\end{observation}

\onslide<+->
\begin{observation}
When $(X,\rho)$ contains all metric spaces and $G$ is an $(n,d)$ graph,
\[ g^{-1}(G,\rho^p) \ge (\log_{d-1}(n/2))^p .\]
\end{observation}
\end{frame}


\begin{frame}{Inembeddability in other spaces}

\onslide<+->
\begin{proposition}
If $(X,\rho)$ a metric and $p\ge 1$, for which there is an $(n,d)$ graphs $(G_n)$,  such that $g^{-1}(G,\rho^p)\le g^{-1}$, then 
\begin{itemize}
\item $c_X(G_n)=\Omega(\log n)$, and
$(G_n)$ do not coarsely embed in $X$.
\end{itemize}
\end{proposition}


\onslide<+->

Theorems of this flavor are  of considerable interest for metric geometry, Banach spaces, and geometric group theory.

Spaces which were studied in this context before.

\begin{itemize}
\item $L_p$ spaces [Matou\v{s}ek]
\item Banach Lattices with cotype [Ozawa]
\item Uniform convex Banach spaces [Kasparov, Yu]
\item CAT(0) spaces [Gromov]
\end{itemize}

\end{frame}


\begin{frame}{Related Previous Results}
\onslide<+->
\begin{theorem}[Matou\v{s}ek]
When $1\le p \le q$, for $(n,d)$ graph $G$,
\[ g(G,(\ell_p,\|\cdot\|_p^p))^{-1} \le g(G,(\ell_q,\|\cdot\|_q^q))^{-1} \le \left( \tfrac{4q}{p} g(G,(\ell_p,\|\cdot\|_p^p))^{-1/p} \right)^q. \]
In particular, any expander family is also expander with respect to $\|\cdot \|_p^p$.
\end{theorem}
\medskip

\onslide<+->

\begin{theorem}[Ozawa, Naor-Rabani]
For any Banach lattice $(X,\|\cdot \|)$ there exist $\alpha:[0,\infty)\to [0,\infty)$
such that
\[ g(G,\|\cdot\|)^{-1} \le \alpha(g(G)^{-1}) \]
\end{theorem}

\end{frame}

\begin{frame}{$B$-Convex Spaces}
\begin{definition}[$B$-convex Banach spaces]
Do not contain $(1+\epsilon)$ distortion of $\ell_1^r$, for some $\epsilon>0$, and $r\ge 2$
\end{definition}

\begin{theorem}[Lafforgue]
If $X$ is $B$ convex banach space then there exist $d\in \mathbb N$, $g^{-1}\ge 1$, and a family of $d$-regular graphs
$(G_n)$, with $g(G_n,\|\cdot\|_X^2)^{-1}\le g^{-1}$.
\end{theorem}

The proof is algebraic using strong form of Property (T).
 \end{frame}

\begin{frame}{Natural questions}

\onslide<+->
\begin{question}
Under what conditions on $(X,\rho)$, there exists $\alpha:[1,\infty)\to [1,\infty)$ such that
$\forall G$,  $g(G,\rho)^{-1} \le \alpha(g(G)^{-1})$ ?
\end{question}

I.e., when does ``expander is expander" ?

\onslide<+-> 
\begin{question}
Under what conditions on $(X,\rho)$ there exist $d\in \mathbb N$,  $g^{-1}\ge 1$, and a family of $d$-regular graphs
$(G_n)$, with $g(G_n,\rho)^{-1}\le g^{-1}$ ?
\end{question} 

I.e., when do expanders exist?

\medskip

\begin{itemize}[<+->]
\item Those questions are quite general. We don't have satisfactory answers.

\item We concentrate on the second question --- building a ``toolbox" for non-linear spectral gaps.

\item Specifically, tools that are instrumental for constructing RVW expanders. 
\end{itemize}
\end{frame}

\begin{frame}{Results}

\begin{itemize}[<+->]
\item A clean, ``abstract" and general analysis of the zigzag product.
\item Geometric condition for the decay of the PI const of powers of graph.
\item A high degree but non-trivial expander graph for $B$-convex spaces in the form of a quotient of a ``noisy Hamming cube".
\item Ball's extension theorem for Lipschitz function into CAT(0) spaces.
\item A unification of non-positive curvature  and uniform convexity via barycentric spaces.
\item RVW expanders for super-reflexive spaces.
\end{itemize}

%
%\onslide<+->
%
%We believe that RVW expanders can be obtained for more general  and other classes of spaces.

\end{frame}

\section{RVW Expanders \& Zig-Zag Product}

\begin{frame}{Review of RVW Expanders  I}

\begin{itemize}[<+->]
\item Reingold, Vadhan , \& Wigderson [RVW] introduced the zigzag product, and used it to construct expanders.
\item RVW expanders were the first explicit expanders  using ``purely" combinatorial methods.
\item Since then, they were used in many applications, some of them impossible using previous methods.
\begin{itemize}
\item Edge expansion beyond the spectral gap. [CRVW]
\item Symmetric log space = log space [Reingold].
\item Expander codes 
\end{itemize}
\item This current work is yet another application of this paradigm.
\end{itemize}
\end{frame}

\begin{frame}{Review of RVW Expanders II}
Ingredients:
\begin{itemize}[<+->]
\item  \alert{Base graph}: A $(D^{2t},D)$ graph $H$, with  $g^{-1}(H) \le \sqrt{t}$.
\item  \alert{Power operation}: Given $(n,d)$ graph $G^s$ is an $(n,d^s)$, with 
$ 1- g(G^s)= (1- g(G))^s$, and hence for odd $s$,
\[ g(G^s)^{-1}  % = (1 -(1-g(G))^s)^{-1}
\le 2 \cdot \max\{1,g(G)^{-1}/s\} \]

Powering: {\color{blue}Improves the PI const $g^{-1}$}, while {\color{red} increasing the degree}.

\item \alert{Zig-zag product}: 
If $G$ is $(n_1,d_1)$-graph, and $H$ is $(n_2,d_2)$ graph, and $n_2=d_1$, then
$G \oz H$ is $(n_1 n_2, d_2^2)$ graph with 
\[ g(G\oz H)^{-1}\le g(G)^{-1} \cdot g(H)^{-2} \]

{Zig-zag: \color{blue} Increase the size, decrease the degree, \color{red}controlled increase of $g^{-1}$}.

\end{itemize}
\end{frame}

\begin{frame}{Review of RVW Expanders III}
\begin{itemize}[<+->]
\item The expander family $(G_i)_i$ is constructed as follows:
\begin{itemize}
\item $G_0=H^2$
\item $G_{i+1}=G_i^t \oz H$
\end{itemize}

\item By induction: $G_i$ is $(D^{2ti},D^2)$ graph with 
\[ g(G_{i+1})^{-1} \le 2 \max\{1, g(G_i)^{-1}/t\} \cdot g(H)^{-2} .\]

\item
Hence, if $g(H)^{-1} \le \sqrt{t/2}$, we have
\[\forall i, \ g(G_i) \le O_t(1) 	 \]
\end{itemize}
\end{frame}

\begin{frame}{Zig-Zag Product [RVW]}
Given:
\medskip

\includegraphics<+>[page=1]{zigzag-op1}
\includegraphics<+->[page=2]{zigzag-op1}

\onslide<+->
\medskip

Define:
$K=G\oz H$ as follows:
\end{frame}

\begin{frame}{Zig-Zag Product [RVW]}
\begin{center}
\includegraphics<+>[page=1 ,scale=0.75]{zigzag-op2}
\includegraphics<+>[page=2 ,scale=0.75]{zigzag-op2}
\includegraphics<+>[page=3 ,scale=0.75]{zigzag-op2}
\includegraphics<+>[page=4 ,scale=0.75]{zigzag-op2}
\includegraphics<+>[page=5 ,scale=0.75]{zigzag-op2}
\includegraphics<+>[page=6 ,scale=0.75]{zigzag-op2}
\includegraphics<+>[page=7 ,scale=0.75]{zigzag-op2}
\includegraphics<+>[page=8 ,scale=0.75]{zigzag-op2}
\includegraphics<+>[page=9 ,scale=0.75]{zigzag-op2}
\includegraphics<+>[page=10,scale=0.75]{zigzag-op2}
\includegraphics<+>[page=11,scale=0.75]{zigzag-op2}
\includegraphics<+>[page=12,scale=0.75]{zigzag-op2}
\includegraphics<+>[page=13,scale=0.75]{zigzag-op2}
\includegraphics<+>[page=14,scale=0.75]{zigzag-op2}
\includegraphics<+>[page=15,scale=0.75]{zigzag-op2}
\includegraphics<+->[page=16,scale=0.75]{zigzag-op2}
\onslide<+->

$G \oz H$ is $(n_1d_1, d_2^2)=(n_1n_2,d_2^2)$ graph.
\end{center}
\end{frame}

\begin{frame}{PI Const of the Zig-Zag Product}

\onslide<+->
[RVW,RTV] showed\onslide<4->{\alert{(*)}} that $g(G\oz H)^{-1} \le g(G)^{-1}\cdot g(H)^{-2}$.

\onslide<+->
\begin{theorem}[{[MN]} \onslide<4->{\alert{*}}]
For every symmetric kernel $\rho$,
\[ g(G\oz H,\rho)^{-1} \le g(G,\rho)^{-1}\cdot g(H,\rho)^{-2}. \]
\end{theorem}

\medskip
\onslide<+->
Of interest even if one only cares about ``classical" spectral gap.

\end{frame}


\begin{frame}{Double Cover (DC) PI Constant}
\begin{overprint}
\begin{minipage}[b]{2.7in}
\begin{overprint}
\onslide*<1-3>{
\begin{definition}[PI Constant]
\[ g(G,\rho)^{-1} = \sup_{\substack{f:V\to X\\ f \ne \text{const}}} \frac{\Avg_{x,y\in V} \rho(f(x),f(y)) }{\Avg_{(x,y)\in E} \rho(f(x),f(y))}. \]
\end{definition}}
\onslide*<4->{
\begin{definition}[DC PI Constant]
\[ \tilde g(G,\rho)^{-1} = \sup_{\substack{f,h:V\to X\\ ~ }} \frac{\Avg_{x,y\in V} \rho(f(x),h(y)) }{\Avg_{(x,y)\in E} \rho(f(x),h(y))}. \]
\end{definition}}
\end{overprint}

\onslide<6->{
\vskip -0.3in
\begin{claim}
\begin{itemize}
\item $\forall G,\rho$, \quad $g(G,\rho)^{-1} \le \tilde g(G,\rho)^{-1}$
\item $\tilde g(G)=1-\lambda(G)$, where $\lambda(G)=\max_{i\ge 2} |\lambda_i(G)|$
\end{itemize}
\end{claim}}
\end{minipage}
\;
\begin{minipage}[t]{1.9in}
\includegraphics<1 >[page=1,width=2in]{dc-pi}
\includegraphics<2 >[page=2,width=2in]{dc-pi}
\includegraphics<3 >[page=3,width=2in]{dc-pi}
\includegraphics<4 >[page=4,width=2in]{dc-pi}
\includegraphics<5->[page=5,width=2in]{dc-pi}
\end{minipage}
\end{overprint}

\onslide<7->
\vskip -0.2in
\begin{theorem}[{[MN]} correct formulation]
$\forall$ symmetric kernel $\rho$, \quad
\( \tilde g(G\oz H,\rho)^{-1} \le \tilde g(G,\rho)^{-1}\cdot \tilde g(H,\rho)^{-2} 
%\onslide*<+>{, \qquad g(G\oz H,\rho)^{-1} \le g(G,\rho)^{-1}\cdot \tilde g(H,\rho)^{-2}}
. \)
\end{theorem}

\onslide<8->
The analysis is simple, but beyond my drawing capabilities. 

We will instead analyze the tensor product $G\otimes H$.
\end{frame}

\begin{frame}{Tensor Product}
Given:
\medskip

\includegraphics[page=1]{tensor-op1}

\medskip

Define:
$K=G\otimes H$ as follows:
\end{frame}

\begin{frame}{Tensor Product}
\begin{center}
\includegraphics<+>[page=1 ,scale=0.70]{tensor-op2}
\includegraphics<+>[page=2 ,scale=0.70]{tensor-op2}
\includegraphics<+>[page=3 ,scale=0.70]{tensor-op2}
\includegraphics<+>[page=4 ,scale=0.70]{tensor-op2}
\includegraphics<+>[page=5 ,scale=0.70]{tensor-op2}
\includegraphics<+>[page=6 ,scale=0.70]{tensor-op2}
\includegraphics<+>[page=7 ,scale=0.70]{tensor-op2}
\includegraphics<+>[page=8 ,scale=0.70]{tensor-op2}
\includegraphics<+>[page=9 ,scale=0.70]{tensor-op2}
\includegraphics<+>[page=10,scale=0.70]{tensor-op2}
\includegraphics<+>[page=11,scale=0.70]{tensor-op2}
\includegraphics<+>[page=12,scale=0.70]{tensor-op2}
\includegraphics<+>[page=13,scale=0.70]{tensor-op2}
\includegraphics<+>[page=14,scale=0.70]{tensor-op2}
\includegraphics<+>[page=15,scale=0.70]{tensor-op2}
\includegraphics<+>[page=16,scale=0.70]{tensor-op2}
\includegraphics<+->[page=17,scale=0.70]{tensor-op2}

\onslide<1->{

\(E(G\otimes H)= \{((u,x),(v,y)):\  (u,v)\in E(G),\ (x,y)\in E(H) \} \)

~\\

\( G \otimes H \text{ is } (n_1 n_2, d_1 d_2) \text{ graph}. \)
}
\end{center}
\end{frame}

\begin{frame}{PI Const of Tensor Product}
Let $A_G$ the adjacency matrix of $G$. 
\begin{claim}[Easy]
$A_{G\otimes H}= A_G\otimes A_H$.
\end{claim}
The eigenvalues of $G\otimes H$ are $\{\lambda(G)_i \lambda(H)_j: i\in [n_1],\, j\in [n_2]\}$.

Therefore $\lambda(G\otimes H)=\max \{\lambda(G),\lambda(H)\}$.

Which means \[\tilde g(G\otimes H)^{-1}= \max\{\tilde g(G)^{-1}, \tilde g(H)^{-1}\}.\]

\onslide<2->{
We will prove
\begin{proposition}
$\forall \rho$,\quad $\onslide*<2>{\tilde g(G\otimes H, \rho)^{-1}\le \tilde g(G,\rho)^{-1} \cdot \tilde g(H,\rho)^{-1}}$
\end{proposition}}
\end{frame}

\begin{frame}{Analysis of Tensor Product}
\begin{center}
\includegraphics<+>[page=1 ,scale=0.60]{tensor-analysis}
\includegraphics<+>[page=2 ,scale=0.60]{tensor-analysis}
\includegraphics<+>[page=3 ,scale=0.60]{tensor-analysis}
\includegraphics<+>[page=4 ,scale=0.60]{tensor-analysis}
\includegraphics<+>[page=5 ,scale=0.60]{tensor-analysis}
\includegraphics<+>[page=6 ,scale=0.60]{tensor-analysis}
\includegraphics<+>[page=7 ,scale=0.60]{tensor-analysis}
\includegraphics<+>[page=8 ,scale=0.60]{tensor-analysis}
\includegraphics<+>[page=9 ,scale=0.60]{tensor-analysis}
\includegraphics<+>[page=10,scale=0.60]{tensor-analysis}
\includegraphics<+>[page=11,scale=0.60]{tensor-analysis}
\includegraphics<+>[page=12,scale=0.60]{tensor-analysis}
\includegraphics<+->[page=13,scale=0.60]{tensor-analysis}
\end{center}

\vskip -0.5in

\begin{align*} 
\uncover<3-> {\Avg_{u,v\in V_H} & \Avg_{x,y\in V_G}  \rho(f(x,u)\,,\,h(y,v)) \\  }
\uncover<10->{  &\le  \tilde g(G,\rho)^{-1} \Avg_{(x,y)\in E_G} \Avg_{u,v\in V_H} \rho(f(x,u)\,,\,h(y,v)) \\}
\uncover<13->{ &\le  \tilde g(G,\rho)^{-1} \cdot \tilde g(H,\rho)^{-1} \Avg_{(x,y)\in E_G} \Avg_{(u,v)\in E_H} \rho(f(x,u)\,,\,h(y,v)) \\}
\uncover<+->{ & = \tilde g(G,\rho)^{-1} \cdot \tilde g(H,\rho)^{-1} \Avg_{((x,u),(y,v))\in E_{G\otimes H}} \rho(f(x,u)\,,\,h(y,v)) }
\end{align*}

\end{frame}

\begin{frame}{Wrapping up the zig-zag product}

\begin{itemize}[<+->]
\item Similar (but slightly more intricate)  analysis  proves the best known bound on $g(G\oz H, \rho)^{-1}$ in surprising generality.

\item The zig-zag  step in RVW expanders works in a general setting.\\
It implies  the existence of  expanders in the most general setting.\\
 {\color{blue}Right?}

\item \alert{Wrong!}

\item From our perspective the zig-zag product is the easiest step in RVW expanders --- almost an ``abstract nonsense".

\item Finding geometric conditions when powering improves the PI constant and finding a base graph is much more challenging.

\item We manage to find non-trivial sufficient conditions.
\alert{But}, we believe that weaker or other conditions are possible.
\end{itemize}
\end{frame}


\section{Decay of the PI Constant}

\begin{frame}{Decay of the PI Const}

\begin{itemize}[<+->]
\item  $\lambda_i(G^t)=\lambda_i(G)^t$, and therefore $\tilde g (G^t)^{-1} \le 2 \max\{1, \tilde g(G)^{-1}/t\}$.

\item We look for a similar effect in more general setting.

\begin{question}[Linear decay of the PI const]
Under what conditions on $\rho$,  $\tilde g(G^t,\rho) \le C\max\{1,\tilde g(G,\rho)/t\}$ ?
\end{question}

\item For RVW-expanders, much weaker decay may suffice.

\begin{question}[Uniform decay of the PI const] 
For which $\rho$, $\forall M\ge 1$,  $\exists t\in\mathbb N$,  $s_0\in[1,\infty)$ such that
$\forall G$, if $\tilde g(G,\rho)\ge s_0$, then $\tilde g (G^t,\rho) \le \tilde g(G,\rho)/M$.
\end{question}

\item Note that some decay is guaranteed.
\begin{proposition} [Non uniform decay of the PI const]
$\forall \rho, G$,\ if $\tilde g(G,\rho)^{-1}<\infty$ then $\lim _{t\to \infty}\tilde g(G^t,\rho)^{-1}=1$
\end{proposition}

\end{itemize}


%
%\onslide<+->
%
%One can ask variants of those question, which may be instrumental to RVW construction, depending on the nature of the base graph.
%
\end{frame}







\begin{frame}{Sufficient condition: Barycentric spaces} 

\onslide<+->

\begin{definition}[$(X,\rho)$ is  $p$-barycentric (w/ const $K$)]
$\forall$ discrete r.v. $U:(\Omega,\mu) \to X$, $\exists \cm(U)$ satisfying
\begin{enumerate}
\item $\forall$ discrete probability space $(\Omega,\mu)$ and  two r.v. $U,V:\Omega\to X$, \\
\(
 \rho(\cm(U),\cm(V)) \le \Avg_\mu \rho(U,V) . 
 \)
\item $\forall$ r.v. $U:(\Omega,\mu)\to X$, and for every $x\in X$, \\
\(
 \rho(\cm(U),x)^p+ K^{-p} \Avg_\mu \rho(U,\cm(U))^p \le \Avg_\mu\rho(U,x)^p.
\)
\end{enumerate}
\end{definition}

\onslide<+->
\begin{example}[Hilbert space is 2-barycentric w/ Const 1]
Take $\cm(U)=\Avg U$, and for every $x\in H$,
\begin{multline*}
 \Avg \|U-x\|^2= \Avg\| (U-\Avg U) - (x-\Avg U)\|^2 \\ =  
\Avg \| U-\Avg U \|^2 \cancel{-2 \langle x -\Avg U,\Avg U -\Avg U \rangle} +\|x -\Avg U\|^2 
\end{multline*}
\end{example}
%
%\begin{example}
%$2$-barycentric spaces with constant $1$ are precisely spaces of non-positive curvature in the sense of Alexandrov (CAT(0)).
%\end{example}
%
%\onslide<+->
%\begin{example}
%$p$-barycentric Banach spaces are precisely the $p$-convex Banach spaces.
%\end{example}

\onslide<+->
The class of barycentric spaces is much richer\ldots
\end{frame}

\begin{frame}{CAT(0): Non-positive curvature}

Geodesic metric spaces with the property:

\medskip
\begin{center}
\includegraphics<+>[page=1]{cat0}
\includegraphics<+>[page=2]{cat0}
\includegraphics<+>[page=3]{cat0}
\includegraphics<+>[page=4]{cat0}
\includegraphics<+>[page=5]{cat0}
\includegraphics<+>[page=6]{cat0}
\includegraphics<+>[page=7]{cat0}
\includegraphics<+>[page=8]{cat0}
\includegraphics<+>[page=9]{cat0}
\includegraphics<+>[page=10]{cat0}
\includegraphics<+>[page=11]{cat0}
\includegraphics<+->[page=12]{cat0}
\end{center}
\end{frame}

\begin{frame}{CAT(0) are 2-barycentric w/ const 1}
\alert{Claim:} \(
 \rho(\cm(U),x)^2+ 1^{-2} \Avg_\mu \rho(U,\cm(U))^2 \le \Avg_\mu\rho(U,x)^2.
\)

\begin{center}
\includegraphics<1>[page=2]{cat0-barycentric}
\includegraphics<2>[page=3]{cat0-barycentric}
\includegraphics<3>[page=4]{cat0-barycentric}
\includegraphics<4>[page=5]{cat0-barycentric}
\includegraphics<5->[page=6]{cat0-barycentric}
\end{center}
\onslide<1->
\begin{multline*} \onslide<4->{\color{red}\rho(\cm(U),x)^2+ 1^{-2} \Avg_\mu \rho(U,\cm(U))^2} {\color{black}\le}  
\onslide<3,4>{\color{green}\|\cm(U)-x\|_2^2+  \Avg_\mu \|U-\cm(U)\|_2^2} \\
=
 \onslide<3>{\color{blue}\Avg_\mu\|U-x\|_2^2}  {\color{black}=} \onslide<3,5->{\color{blue}\Avg_\mu\rho(U,x)^2}.
\end{multline*}

\end{frame}

\begin{frame}{Uniform Convex Banach Spaces}
\begin{center}
\includegraphics<+>[page=1]{uconvex}
\includegraphics<+->[page=2]{uconvex}
\end{center}
\begin{definition}[Uniformly convex Banach space $X$]
\( \forall \e>0\; \exists \delta>0\; \forall x,y\in X,\; \text{ if } \|x\|_X=\|y\|_X=1, 
\text{ then } \|\frac{x+y}{2}\| \le 1-\delta. \)
\end{definition}
\onslide<3->
\begin{theorem}[{[Pisier] $p$-convexity}]
U. Convex space can be renormed s.t.  $\exists 2\le p<\infty$, and $K>0$
\[ \forall x,y\in X,\; \left\|\frac{x+y}{2}\right\|^p + K^{-p} \left\|\frac{x-y}2 \right \|^p \le  \frac{\|x\|^p+\|y\|^p}2 \]
\end{theorem}
\end{frame}

\begin{frame}{Uniform Convex are $p$-barycentric for some $p<\infty$}
\onslide<+->
\begin{theorem}[{[Pisier] $p$-convexity}]
U. Convex space can be renormed s.t.  $\exists 2\le p<\infty$, and $K>0$
\[ \forall x,y\in X,\; \left\|\frac{x+y}{2}\right\|^p + K^{-p} \left\|\frac{x-y}2 \right \|^p \le  \frac{\|x\|^p+\|y\|^p}2 \]
\end{theorem}

\onslide<+->
\begin{corollary}
$\forall$ p.d. $U$,
\( \|\Avg U \|^p + K^{-p} \Avg \|U-\Avg U\|^p \le \Avg \|U\|^p .\)
\end{corollary}

\onslide<+->
This is the 2nd item of $p$-barycentric space
\[
 \rho(\cm(U),x)^p+ K^{-p} \Avg_\mu \rho(U,\cm(U))^p \le \Avg_\mu\rho(U,x)^p,
\]
where $\cm(U)=\Avg U$, and $x=0$ (translation invariance).

\medskip
\onslide<+->
The 1st item follows from the convexity of the norm.
\end{frame}


\begin{frame}{The plan for barycentric spaces}

\begin{enumerate}[<+->]
\item Deduce Pisier's martingale inequality for barycentric spaces.
\item The martingale inequality implies a strong form of Ball's metric Markov cotype inequality.
\item Ball's metric Markov cotype inequality implies the decay of $g^{-1}$:\\ 
Let $\mathcal B_t(G) =\tfrac{1}{t}\sum_{i=0}^{t-1} A_G^i$, then
\[ g^{-1}(\mathcal B_t(G),\rho^p) \le C \max\{1,g^{-1}(G,\rho^p)/t \} .\]
\item 
This decay falls short of power of graph, 
but it is good enough for the RVW construction since the degree is not much larger.

%\item Side note:  A Lipschitz extension theorem from $L_p$ spaces, $p\ge 2$, to CAT(0) spaces.
%Follows from the metric Markov cotype 2 property, and Ball's extension theorem.
\end{enumerate}
\end{frame}

\begin{frame}{Martingales in Barycentric Spaces}

\onslide<1->{
\begin{definition}[Martingale in $p$-barycentric space $(X,\rho)$]
$(M_n)_n$: A seq' of $X$ valued variables of finite prob. space $(\Omega, \mathcal F, \mu)$, for which there is filterization $(\mathcal F_n))_n$
of $\mathcal F$ such that
\[ \cm\left (M_{n+1}| \mathcal F_n \right)=M_n .\]
\end{definition}
}

\onslide<2->{
\begin{example}[It is possible that $\cm(M_{n+2}| \mathcal F_n)\ne M_n$]
\begin{center}
\includegraphics<1-2>[page=1]{martingale-not}
\includegraphics<3>[page=2]{martingale-not}
\includegraphics<4>[page=3]{martingale-not}
\includegraphics<5>[page=4]{martingale-not}
\includegraphics<6>[page=5]{martingale-not}
\includegraphics<7>[page=6]{martingale-not}
\includegraphics<8>[page=7]{martingale-not}
\includegraphics<9>[page=8]{martingale-not}
\includegraphics<10>[page=9]{martingale-not}
\includegraphics<11>[page=10]{martingale-not}
\includegraphics<12>[page=11]{martingale-not}
\includegraphics<13>[page=12]{martingale-not}
\end{center}
\end{example}
}
\end{frame}

\begin{frame}{Pisier's Martingale Cotype Inequality}

\onslide<+->
\begin{proposition}
Let $(M_n)_{n=0}^m$ be a martingale in $p$-barycentric space $(X,\rho)$, then
\[ \sum_{i=0}^{m-1} \Avg \rho(M_i, M_{i+1})^p \le K^p \Avg \rho(M_m,M_0)^p .\]
\end{proposition}
\onslide<+->
\begin{proof}
Recall that \(
 \rho(\cm(U),x)^p+ K^{-p} \Avg_\mu \rho(U,\cm(U))^p \le \Avg_\mu\rho(U,x)^p.
\)

Take $U=[M_{i+1}|\mathcal F_i]$, and $x=M_0$, and obtain
\[ \rho(M_i,M_0)^p+ K^{-p} \Avg_\mu [\rho(M_{i+1},M_i)^p|\mathcal F_i] \le \Avg_\mu[\rho(M_{i+1},M_0)^p|\mathcal F_i]. \]

Summing over $i$ and taking expectation over $\mathcal F$, we conclude observing the telescoping sum.
\end{proof}
\end{frame}

\begin{frame}{Markov Cotype}

\onslide<+->
\begin{definition}[(metric) Markov cotype {[Ball]}]
A metric space $(X,\rho)$ has Markov cotype $p$ w/ const $K>0$ 

 $\forall n\times n \text{ symmetric stochastic matrix }A,$
 \vskip -0.2in
\begin{multline*}
\  \forall x_1,\ldots x_n\in X, \ \forall \alpha\in (0,1),\ \exists y_1,\ldots, y_n\in X \\
 \frac{\alpha}{1-\alpha} \sum_{i,j} A_{ij}\rho(y_i,y_j)^p+ 2\sum_i\rho(y_i,x_i)^p \le K^p\sum_{i,j}
 \mathcal G_\alpha(A)_{i,j} \rho(x_i,x_j)^p, 
 \end{multline*}
 \vskip -0.2in
 where  $\mathcal G_\alpha(A)=(1-\alpha)(I-\alpha A)^{-1}=(1-\alpha) \sum_{\ell=0}^\infty (\alpha A)^\ell $.
\end{definition}

\begin{overprint}
\begin{itemize}
\onslide*<2-3>{\item
{[Ball]}: A Lipschitz mapping $f:S\to X$ can be extended
to $Y\supset S$ provided that  $X$ has \alert{Markov cotype 2} and  $Y$ has a property called \alert{Markov type 2}.
}
\onslide*<3>{\item We use Markov cotype 2 to prove the decay of the PI const under $\mathcal G_\alpha$ operation.}

\item<4-> Writing $\alpha\approx 1 -1/t$, \ 
\( \mathcal G_\alpha(A)\approx \sum_{i={0}}^\infty e^{-i} A^{it} \frac{1}{t} \sum_{j=0}^{t-1} A^j .\)

\item<5-> I don't see how to represent (or approx.) $\mathcal G_\alpha(A)$ as low deg. graph.


\item<6->  If $A=A_G$, and $G$ is $(n,d)$ graph,  $\frac{1}{t} \sum_{j=0}^{t-1} A^j $ can be represented by
$(n,td^t)$ graph, which is good enough. 
\end{itemize}
\end{overprint}
\end{frame}

\begin{frame}{Variant of the Markov Cotype}
\onslide<+->
\begin{definition}[Strong Markov cotype]
A metric space $(X,\rho)$ has strong Markov cotype $p$ w/ const $K>0$ 

 $\forall n\times n \text{ symmetric stochastic matrix }A,$
 \vskip -0.2in
\begin{multline*}
  \forall x_1,\ldots x_n\in X, \ \forall t\in \mathbb N,\ \exists y_1,\ldots, y_n\in X \\
 t \sum_{i,j} A_{ij}\rho(y_i,y_j)^p+ 2\sum_i\rho(y_i,x_i)^p \le K^p\sum_{i,j}
 \mathcal B_t(A)_{i,j} \rho(x_i,x_j)^p, 
 \end{multline*}
 \vskip -0.2in
 where  $\mathcal B_t(A)= \frac 1t \sum_{i=0}^{t-1} A^i $.
\end{definition}

Strengths:
\begin{itemize}[<+->]
\item Can be proved from martingale cotype.  Ball proves Markov cotype only for uniformly convex spaces.
\item implies Markov cotype.
\item Can be represented with relatively low degree graph:\\
\( \mathcal B_t(G)= \bigcup_{i=0}^{t-1} d^{t-1-i} \times G^i . \)
\end{itemize}
\end{frame}

\begin{frame}{Proof of `Strong Markov Cotype` in Barycentric Space}
\begin{itemize}[<+->]
\item Given: $n\times n$ matrix $A$, $t\in \mathbb N$, and $x_1,\ldots, x_n\in X$.
\item Fix $\ell\in [n]$.
\item Define  filterization $(\mathcal F_s)_{s=1}^t$.  $\mathcal F_s=[n]^s$,
\[ \Pr\bigl[(i_0=\ell,i_1,\ldots, i_s)\bigr ]=\prod_{i=1}^s A_{i_{j-1} i_j}. \]
\item Define martingale $(M_s)_{s=0}^t$. $M_t(\ell,i_1,\ldots, i_t)=x_{i_t}$,
\[ M_s(i_0=\ell,i_1,\ldots, i_s)= \cm \Bigl( \sum_j A_{i_s,j}M^k_{s+1}(i_0=\ell,i_1,\ldots, i_s,j) \Bigr ) .\]
\item Apply the martingale cotype ineq., and average over all $\ell$.
\item Define $y_\ell=\cm\bigl (\tfrac{1}{t} \sum_{k=1}^t M_k(\ldots, \ell) \bigr)$. 
\item Do 2 pages of (straightforward) calculations.
\end{itemize}
\end{frame}

\begin{frame}{The Decay of of the PI Const}
\onslide<+->
\begin{proposition}
If $(X,\rho)$ has Markov cotype $p$, then for any graph $G$, \\
\(g(\mathcal B_t(G),\rho^p)^{-1} \le O(1) \cdot \max\{1,g(G,\rho^p)^{-1}/t\} .\)
\end{proposition}
\begin{proof}
Rewriting the conclusion:

$\forall n\times n$ matrix $A$, $\forall x_1,\ldots,x_n\in X$, and $\forall t\in \mathbb N$,
$\exists y_1,\ldots, y_n\in X$ s.t.
\[ \frac{\sum_{i,j} \mathcal B_t(A)_{i,j} \rho(x_i,x_j)^p}{n^{-1} \sum_{i,j} \rho(x_i,x_j)^p} \ge
\Omega(1) \cdot \min \biggl \{1, \frac{ t\cdot \sum_{i,j}  A_{i,j} \rho(y_i,y_j)^p}{n^{-1} \sum_{i,j} \rho(y_i,y_j)^p} \biggr \}
\]
The numerators are close to the Markov cotype inequality ---

\smallskip
\begin{center}
\(  K^p\sum_{i,j}
 \mathcal B_t(A)_{i,j} \rho(x_i,x_j)^p \ge 
t \text{$\sum_{i,j}$} A_{ij}\rho(y_i,y_j)^p+ 2\sum_i\rho(y_i,x_i)^p .\)
\end{center}
\smallskip

The denominators discrepancy is mitigated by $\sum_i \rho(y_i,x_i)^p$.
\end{proof}

\end{frame}

\begin{frame}{Wrapping up the decay of the PI Const}

\begin{itemize}[<+->]
\item In  CAT(0) and uniformly convex  the PI const decays linearly.
\item The proof used a series of reductions: 
\begin{multline*}
 \text{CAT(0)/U. Convex} \longrightarrow \text{Barycentric space} \longrightarrow \text{Martingale cotype}\\
 \longrightarrow \text{strong Markov cotype} \longrightarrow \text{Decay of PI const.}
 \end{multline*}
\item In the paper: the decay of the DC PI const --- $\tilde g(\mathcal B_t(G),\rho^p)$.
\item Side benefit: Lipschitz extensions of functions from subset of $L_p$ or series/parallel graphs, into CAT(0) spaces.
\item It will be more satisfying to have the decay for $G^t$,  instead of $\mathcal B_t(G)$.
\item For RVW expanders both the hypothesis (barycentric) and the conclusion (\alert{linear} decay) may be too strong.
\end{itemize}
\end{frame}


\section{Base graph in $B$-convex Banach spaces}

\begin{frame}{Base Graph: An Introduction}
The ``base graph" for $(X,\rho)$:
\begin{itemize}
\item A family of $(N,N^{1/ \log^{\delta}N})$ graphs $G_N$, where $\delta=\delta(X)>0$. 
\item For each $G_N$, $g(G_N,\rho^p)^{-1}=O_{\rho,p}(1)$.
\item The point: Assuming uniform decay, we can choose $N$ large enough, such that the degree is small enough compared to the size, allowing just high enough powers.
\item The family of graphs will be based on (but different from) the Hamming cubes.
\end{itemize}
\end{frame}


\begin{frame}{Hamming Cube in $L_2$}
Let $C_n=(\{\pm 1\}^n,\; \{(x,y)\in X_n\times C_n:\; \|x-y\|_1=2\})$

\begin{theorem}[Enflo]
\[ \forall f:C_n \to L_2,\  \Avg_{x,y\in C_n}\|f(x)-f(y)\|_2^2 \le O({n}) \!\!\Avg_{(x,y)\in E(C_n)} \!\! \|f(x)- f(y)\|_2^2 .\]
\end{theorem}
\begin{proof}[Khot-Naor]
For $A\subset [n]$, let $W_A:\{\pm 1\}^n \to \{\pm 1\}$, $W_A(x)=\prod_{i\in A} x_i$.
Then
\[ f=\sum_A \hat f(A) W_A, \qquad \text{where}, \qquad \hat f(A)= \Avg_{x\in C_n} f(x) W_A(x)  \in L_2. \]
As we have seen, $\Avg_{x,y\in C_n}\|f(x)-f(y)\|_2^2\sim \Avg_x \|f(x)\|_2^2$, whereas
\( \Avg_{(x,y)\in E(C_n)} \!\!\! \|f(x)- f(y)\|_2^2\sim (1-\lambda_2 \)
\end{proof}
\end{frame}

\begin{frame}{Improving the PI for the Cube [Khot-Naor]}
\begin{itemize}[<+->]
\item Crucial observation I:
If $f:C_n \to L_2$ is supported only on the eigenfunctions of low eigenvalues, then the PI const improves.

\item Crucial Observation II:
This happens when $f$ is const on $x+C^\bot$, and $C$ is code of large weight $w$. In this case, if $|A|<w$, then $\hat f(A)=0$.

\item Hence: Take a linear code $D\subset \{\pm 1\}^n$, $\dim(D)=\Omega(n)$, and $w(D)=\Omega(n)$.
Then $C_n/D^\bot$ is a graph with $2^{\Omega(n)}$ vertices, and whose (effective) degree is still $n$,
and

\begin{multline*} \forall f:C_n/D^\bot \to L_2,\\  
\Avg_{x,y\in C_n}\|f(x)-f(y)\|_2^2 \le O({1}) \!\!\Avg_{(x,y)\in E(C_n)} \!\! \|f(x)- f(y)\|_2^2 .
\end{multline*}

\end{itemize}
\end{frame}


\begin{frame}{The Laplacian}

\end{frame}

\begin{frame}{Bonami-Beckner Operator}

\end{frame}

\begin{frame}{Bounding the Laplacian via BB operator}

\end{frame}

\begin{frame}{Bounding BB on $K$-convex spaces}

\end{frame}

\begin{frame}{Near Miss?}

\end{frame}

\begin{frame}{``Noisy Cube''}

\end{frame}


\section{Discussion}

\begin{frame}{Concrete Construction}

Since uniformly convex Banach spaces are B-convex, we have
all three ingredients for RVW expanders in uniformly convex Banach spaces.

\onslide<+->
\begin{theorem}
There are RVW expanders w.r.t uniformly convex Banach spaces.
\end{theorem}

\onslide<+->
For B-convex and CAT(0) spaces we have 2-out-3 ingredients. 
\begin{question}
\begin{itemize}
\item Is there a base graph for CAT(0) spaces?
\item Is there a uniform decay of the PI const in B-convex spaces?
\end{itemize}
\end{question}

\end{frame}

\begin{frame}{Open Problems}
\onslide<+->
We already asked:
\begin{question}[Q1]
Under what conditions on $(X,\rho)$, there exists $\alpha:[1,\infty)\to [1,\infty)$ such that
$\forall G$,  $g(G,\rho)^{-1} \le \alpha(g(G)^{-1})$ ?
\end{question}

\begin{question}[Q2]
Under what conditions on $(X,\rho)$ there exist $d\in \mathbb N$,  $g^{-1}\ge 1$, and a family of $d$-regular graphs
$(G_n)$, with $g(G_n,\rho)^{-1}\le g^{-1}$ ?
\end{question} 

\onslide<+->
\begin{question}
If and when we have Q2 $\land$ $\lnot$ Q1:\\
Can this stronger expansion property be used in other contexts?
\end{question}

\end{frame}


\begin{frame}{Agenda}

\onslide<+->
Beyond those concrete questions. An agenda emerges:

\begin{quote}
``Calculus of non-linear spectral gaps"
\end{quote}

We feel that this subject may be a ground for a rich theory.

\end{frame}

\begin{frame}{My Pet Problem}
\onslide<+->{
(Far away) Goal:
\begin{question}[Dichotomy]
For any metric space $(X,\rho)$.
Either 
\begin{itemize}
\item $c_X(M)=1$, $\forall$ finite $M$
\item  $\exists (M_n)_n$, such that $|M_n|=n$, and $c_X(M_n)=\Omega(\log n)$.
\end{itemize}
\end{question}}
\onslide<+->{
We do know:
\begin{theorem}[Weaker Dichotomy {[M.--Naor]}]
For any metric space $(X,\rho)$.
Either 
\begin{itemize}
\item $c_X(M)=1$, $\forall$ finite $M$
\item  $\exists \beta>0$, $\exists (M_n)_n$, such that $|M_n|=n$, and $c_X(M_n)=\Omega(\log^\beta n)$.
\end{itemize}
\end{theorem}}

\end{frame}

\begin{frame}{The Dichotomy \& Expanders}
\begin{itemize}[<+->]
\item An  approach for proving the dichotomy: Given $(X,\rho)$, either
\begin{enumerate}
\item $X$ contains all finite spaces (``cotype $\infty$'').
\item $X$ does not (``finite cotype''), and then contains a base graph.
\end{enumerate}
\item Then we also need general uniform  decay of the PI const .
\item If those conditions are met: On the 2nd case we can construct an expander.
\item This approach fails for general metric spaces. Example:  $\sqrt{\ell_\infty}$.
\item It may be true for normed spaces. May be connected to open questions about metric cotype.
\end{itemize}

\end{frame}

\begin{frame}
\vskip 0.7in

\begin{center}

{\Huge Thank You}
\end{center}
\end{frame}

\end{document}


